\documentclass[11pt]{article}
\usepackage[utf8]{inputenc}
\usepackage{amsmath,amssymb}
\usepackage{graphicx}
\usepackage{booktabs}
\usepackage{hyperref}
\usepackage[margin=1in]{geometry}
\usepackage{natbib}
\bibliographystyle{unsrtnat}

\title{Spatio-Temporal Heterogeneity of Tumour-Residing Mature Regulatory Dendritic Cells and Their Role in Anti-PD-L1 Immunotherapy}
\author{}
\date{}

\begin{document}
\maketitle

\begin{abstract}
Mature regulatory dendritic cells (mRegDCs; LAMP3$^{+}$/CCR7$^{+}$/PD-L1$^{+}$) are identified across human cancers, yet whether they uniformly migrate to draining lymph nodes or persist within tumours remains unclear. Using photoconvertible Kaede mice bearing syngeneic tumours, we demonstrate that while mRegDCs are the dominant DC population emigrating to draining lymph nodes, a subset remains tumour-resident despite expressing CCR7. Tumour-retained mRegDCs undergo progressive transcriptional changes with increasing dwell-time, showing reduced antigen presentation capacity and diminished pro-inflammatory gene expression. In both human and murine solid tumours, mRegDCs spatially co-localize with PD-1$^{+}$CD8$^{+}$ T cells. Anti-PD-L1 treatment skews tumour mRegDCs toward a state enriched in lymphocyte stimulatory molecules, notably OX40L, and OX40L-expressing mRegDCs augment anti-tumour cytolytic CD8$^{+}$ T cell activity. Analysis of 4,045 TCGA tumours and 1,853 METABRIC breast cancers confirms that a mRegDC gene signature associates with improved patient survival across melanoma, breast, colorectal, and lung cancers. These findings reveal that tumour mRegDCs are phenotypically heterogeneous, temporally dynamic, and functionally reshaped by checkpoint blockade to support anti-tumour immunity.
\end{abstract}

\section{Introduction}

Dendritic cells (DCs) are central orchestrators of anti-tumour immunity, bridging innate and adaptive immune responses through antigen uptake, processing, and presentation to T cells \citep{wculek2020dendritic,gardner2020dendritic}. Within the tumour microenvironment (TME), a transcriptionally distinct DC population termed mature regulatory DCs (mRegDCs)---characterized by co-expression of LAMP3, CCR7, and PD-L1---has been consistently identified across multiple human cancer types through single-cell RNA sequencing \citep{maier2020systematic,cheng2021multi}. However, the functional significance of mRegDCs in anti-tumour immunity remains debated, with evidence supporting both immunostimulatory and immunosuppressive roles.

The prevailing model assumes that CCR7 expression on mRegDCs drives their migration from the tumour to draining lymph nodes (dLN), where they activate T cell responses. This assumption has led to mRegDCs being classified primarily as migratory DCs. However, CCR7 expression does not guarantee migration, and whether a population of mRegDCs remains tumour-resident has not been directly tested. If such cells exist, their phenotype and function---shaped by prolonged exposure to the immunosuppressive TME---could differ substantially from their migratory counterparts.

Furthermore, checkpoint immunotherapy with anti-PD-L1 antibodies has transformed cancer treatment, yet the specific effects on mRegDC states and their downstream interactions with CD8$^{+}$ T cells are poorly characterized. Given that mRegDCs express high levels of PD-L1, they represent a direct cellular target of anti-PD-L1 therapy beyond the canonical PD-1/PD-L1 axis on T cells and tumour cells.

Here we use photoconvertible Kaede transgenic mice to directly track DC migration dynamics, combined with single-cell transcriptomics, functional assays, and analysis of large human cancer cohorts, to dissect the heterogeneity, temporal dynamics, and therapeutic responsiveness of tumour-residing mRegDCs.

\section{Results}

\subsection{mRegDCs include a tumour-resident subset despite CCR7 expression}

To directly track DC migration from tumours to draining lymph nodes, we established subcutaneous MC38, MC38-Ova, and CT26 tumours in Kaede transgenic mice (C57BL/6 and BALB/c backgrounds). At day 13, tumour photoconversion labelled all resident cells with Kaede-red fluorescence, enabling identification of tumour emigrants in the dLN as Kaede-red$^{+}$ cells at 5h, 24h, 48h, and 72h post-conversion (Figure~2A--B).

Flow cytometric analysis confirmed that mRegDCs were the dominant DC population arriving in the dLN, with Kaede-red$^{+}$ mRegDCs detectable as early as 5h and peaking at 48h post-conversion. However, CCR7$^{+}$ mRegDCs also remained present in the tumour at all time points examined, demonstrating that a subset of mRegDCs persists in the TME despite expressing the canonical migration receptor (Figure~2). The mRegDC-to-cDC ratio in tumours increased over time in both MC38 and CT26 models (one-way ANOVA with Sidak's correction; Extended Data Figure~3C), suggesting progressive accumulation of tumour-retained mRegDCs.

\subsection{Tumour mRegDCs are transcriptionally heterogeneous and change with dwell-time}

Single-cell RNA-seq of FACS-sorted CD45$^{+}$ cells from MC38-Ova tumours at 48h post-photoconversion revealed multiple transcriptionally distinct mRegDC clusters (Figure~1B--D). These could be ordered along a temporal trajectory corresponding to increasing tumour dwell-time:

\begin{itemize}
    \item \textbf{mRegDC\_1} (recent arrivals): high expression of antigen presentation genes (MHC-II, Cd74) and pro-inflammatory cytokines
    \item \textbf{mRegDC\_2} (intermediate): transitional phenotype
    \item \textbf{mRegDC\_3} (prolonged residence): progressive reduction in MHC-II transcripts, Cd74, and antigen processing pathways (GO:0002468), with maintained or increased co-inhibitory molecules including PD-L1
\end{itemize}

Gene signature score comparisons between clusters confirmed the progressive functional decline (Wilcoxon rank-sum test; Figure~3A--C). PAGA analysis and RNA velocity supported a maturation trajectory from recently arrived to long-term resident mRegDCs within the tumour (Extended Data Figure~3).

Tumour mRegDCs were also transcriptionally distinct from their dLN counterparts, which showed higher expression of lymphocyte co-stimulation pathways (GO:0031294) by GSEA and more uniform MHC-II expression (Extended Data Figure~4). These data collectively demonstrate that the tumour microenvironment progressively reshapes mRegDC phenotype over time, reducing their capacity for antigen presentation and T cell stimulation.

\subsection{Anti-PD-L1 treatment skews mRegDCs toward a stimulatory state}

Mice bearing MC38 tumours were treated with anti-PD-L1 (clone 80, 200~$\mu$g IP on days 7, 10, and 13) or isotype control ($n = 5$ per group). Anti-PD-L1 treatment reduced tumour growth and induced substantial transcriptional changes in tumour mRegDCs (Figure~3, Extended Data Figure~6).

GSEA revealed significant upregulation of hallmark IFN-$\gamma$ response, hallmark inflammatory response, and KEGG antigen processing/presentation pathways in mRegDCs from anti-PD-L1-treated tumours (adjusted $p < 0.05$; Extended Data Figure~6). Key molecules upregulated included OX40L (Tnfsf4), a potent T cell co-stimulatory ligand, alongside maintained high CD80 expression and increased PVR (TIGIT ligand).

RNA velocity analysis revealed that in isotype-treated tumours, mRegDCs progressed toward an exhausted, immunosuppressive state, whereas anti-PD-L1 treatment redirected this trajectory toward a lymphocyte-stimulatory state enriched in co-stimulatory molecules (Extended Data Figure~6C).

\subsection{mRegDCs co-localize with CD8$^{+}$ T cells and support anti-tumour immunity}

Multiplex immunofluorescence demonstrated spatial co-localization of mRegDCs with PD-1$^{+}$CD8$^{+}$ T cells in both human solid tumours and murine MC38-Ova tumours (Figure~4). This spatial proximity was supported by bulk transcriptomic analysis: deconvolution of 4,045 TCGA tumour transcriptomes revealed strong positive correlations between mRegDC and CD8$^{+}$ T cell proportions across colorectal ($n = 521$), breast ($n = 1,226$), and melanoma ($n = 473$) cancers (Pearson correlation, $p < 0.05$; Figure~4A). The mRegDC gene signature was also positively correlated with a CD8$^{+}$ effector T cell signature across these cancer types (Figure~4B).

Functional experiments using OX40L reporter mice confirmed that tumour mRegDCs express OX40L. DC-specific OX40L deletion (CD11c-Cre $\times$ OX40L$^{fl/fl}$ mice) demonstrated that OX40L$^{+}$ mRegDCs augment cytolytic CD8$^{+}$ T cell activity, supporting a model in which anti-PD-L1 therapy enhances anti-tumour immunity in part by upregulating co-stimulatory molecules on tumour-resident mRegDCs.

\subsection{mRegDC signature associates with improved survival across cancer types}

Kaplan--Meier survival analysis of TCGA cohorts demonstrated that high mRegDC gene signature expression was associated with improved overall survival in skin cutaneous melanoma ($n = 473$), colorectal adenocarcinoma ($n = 521$), breast invasive carcinoma ($n = 1,226$), and lung adenocarcinoma ($n = 598$; log-rank test; Extended Data Figure~1). Subtype-specific analysis in the METABRIC breast cancer cohort ($n = 1,853$) revealed associations across HR$^{+}$ ($n = 1,369$), HER2$^{+}$ ($n = 185$), and triple-negative ($n = 299$) subtypes (Extended Data Figure~1H).

Published single-cell RNA-seq datasets from colorectal cancer (62 patients, GSE178341), breast cancer (29 patients, EGAS00001004809), and melanoma (25 patients, GSE123139) consistently identified mRegDC clusters expressing CCR7, PD-L1, and PD-L2, confirming the generalizability of our murine findings to human cancers.

\section{Discussion}

Our study reveals that tumour-resident mRegDCs are not a uniform population but rather exhibit spatio-temporal heterogeneity shaped by tumour dwell-time and therapeutic intervention. The demonstration that CCR7$^{+}$ mRegDCs persist in tumours challenges the prevailing assumption that CCR7 expression equates to lymph node migration. Instead, a fraction of mRegDCs remains trapped in the immunosuppressive TME, where prolonged residence progressively diminishes their antigen presentation capacity.

The therapeutic implications are substantial. Anti-PD-L1 treatment does not merely block PD-L1 on mRegDCs but actively reprograms their transcriptional state, upregulating co-stimulatory molecules including OX40L that enhance CD8$^{+}$ T cell cytolytic activity. This mechanism operates in addition to the canonical effects of checkpoint blockade on T cell-intrinsic signalling, suggesting that mRegDC reprogramming represents an underappreciated component of anti-PD-L1 efficacy \citep{garris2018successful}.

The consistent association between the mRegDC signature and improved survival across four cancer types in TCGA data supports the view that mRegDCs are net positive contributors to anti-tumour immunity, despite their co-expression of inhibitory molecules. The balance between stimulatory and inhibitory functions appears to depend on the temporal dynamics of tumour residence and can be therapeutically modulated.

\section{Methods}

\subsection{Mouse tumour models}
Subcutaneous syngeneic tumours (MC38, MC38-Ova, CT26; 2.5$\times$10$^5$ cells, left flank) were established in photoconvertible Kaede transgenic mice (C57BL/6 and BALB/c). Anti-PD-L1 (clone 80, 200~$\mu$g IP) or isotype control (NIP228 mouse IgG1) was administered on days 7, 10, and 13. Tumour volume was calculated as $V = 0.5 \times a \times b^2$. Mice were housed at 21$^\circ$C, 55\% humidity, on a 12h light-dark cycle.

\subsection{Photoconversion and tracking}
Tumours were exposed to 405~nm light on day 13, labelling resident cells with Kaede-red. DCs were analyzed by flow cytometry in tumours and dLN at 5h, 24h, 48h, and 72h post-conversion.

\subsection{Single-cell RNA-seq}
CD45$^{+}$ cells were FACS-sorted from MC38-Ova tumours and processed for scRNA-seq. Clustering, PAGA, RNA velocity, and gene set enrichment analyses were performed. Gene signature scores were compared using Wilcoxon rank-sum tests.

\subsection{Human data analysis}
TCGA bulk transcriptomes ($n = 4,045$) and METABRIC ($n = 1,853$) were analyzed for survival associations using Kaplan--Meier analysis with log-rank tests. Published scRNA-seq datasets from CRC (GSE178341), breast cancer (EGAS00001004809), and melanoma (GSE123139) were reanalyzed.

\subsection{Functional experiments}
OX40L reporter (OX40L$^{+}$/Human-CD4) and conditional knockout (CD11c-Cre $\times$ OX40L$^{fl/fl}$) mice were used for functional studies. Multiplex immunofluorescence assessed spatial co-localization.

\bibliography{references}

\end{document}
